%!TEX root=../../note.tex
\subsection{Principle of Strong Induction}
Why Mathematical Induction is not enough? 
\begin{example}
    A sequence $\set{x_n }$ is defined by 
    \begin{equation*}
        x_1 = 4 \qquad x_2 = 68 \qquad x_n = 2 x_{n - 1} + 15x_{n - 2} \qquad \forall n \geq 3
    \end{equation*}

    Prove $\underbrace{x_n}_{\text{closed form}} = 2 \paren{-3}^n + 10 \cdot 5^{n - 1}, \forall n \geq 1$.
\end{example}
\begin{remark}
    This is an example of a sequence defined \textit{recursively} or a recursive relation. 
\end{remark}

\begin{proof}
    Let $P(n)$ be $x_n = 2\paren{-3}^n + 10 \cdot 5^{n - 1}$

    \textbf{Base case: } Prove $P(1)$ is true.  We have
    \begin{equation*}
        x_1 = 4
    \end{equation*}

    by definition. also
    \begin{align*}
        2 \times \paren{3}^1 + 10 \cdot 5^{1 - 1} = 2(-3) + 10 = -6 + 10 = 4 = x_1
    \end{align*}

    So $P(1)$ is true. 

    \textbf{Inductive Step}
    Let $k$ be an arbitrary natural number. 

    Assume $P(k)$ is true. 

    That is 
    \begin{equation*}
        x_k = 2\paren{-3}^k + 10 \cdot 5^{k - 1}. 
    \end{equation*}
    
    Prove $P(k + 1)$. That is 
    \begin{align*}
        x_{k + 1} = 2\paren{-3}^{k + 1} + 10 \cdot 5^k. 
    \end{align*}
\end{proof}

We cannot continue with this. Since we are unable to prove that 
\begin{equation*}
    P(k) \implies P(k + 1)
\end{equation*}

since we only know $P(k)$ is true. If we want or prove, we also need to know that $P(k - 1)$ is true. 

\begin{definition}
    [Principle of Strong Induction]
    Let $P(n)$ be an open sentence that depends on $n \in \N$.  

    If statement 1 and 2 are both true 
    \begin{enumerate}
        \item P(1)
        \item For all $k \in \N$, if 
        \begin{equation*}
            P(1) \wedge P(2) \wedge \cdots \wedge P(k)
        \end{equation*}
        is true, then $P(k + 1)$ is true. 
    \end{enumerate}

    then $P(k), \forall k \in \N$. 
\end{definition}

\subsubsection*{Prove by Strong Induction}
Prove $\forall n \in \N, n \geq b, P(n)$ 
\begin{itemize}
    \item Prove $P(b) \wedge P(b + 1) \wedge \cdots \wedge P(B), \forall B \geq b$. 
    \item Inductive step: Prove for all integers $k \geq B$. 
\end{itemize}

Go back to the bloody example. 
\begin{example}
    A sequence $\set{x_n }$ is defined by 
    \begin{equation*}
        x_1 = 4 \qquad x_2 = 68 \qquad x_n = 2 x_{n - 1} + 15x_{n - 2} \qquad \forall n \geq 3
    \end{equation*}

    Prove $\underbrace{x_n}_{\text{closed form}} = 2 \paren{-3}^n + 10 \cdot 5^{n - 1}, \forall n \geq 1$.
\end{example}

\begin{proof}
   Proof by induction on $n$.  

   \textbf{Base Case: }Prove $P(1)$
   By definition, 
    \begin{align*}
        2 \times \paren{-3}^1 + 10 \cdot 5^{1 - 1} = 2(-3) + 10 = -6 + 10 = 4 = x_1 \\
        2 \times \paren{3}^2 + 10 \cdot 5^{2 - 1} = 18 + 50 = 68 = x_2
    \end{align*}

    Therefore, $P(1)$ is true. 

    \textbf{Inductive step }:Let $k$ be an arbitrary integer with $k \geq 2$

    Assume $P(i)$ is true for all integer $i$ with $1 \leq i \leq k$. 

    That is assume 
    \begin{equation*}
        x_i = 2\paren{-3}^i + 10 \cdot 5^{i - 1}, \forall i \in \set{1, 2, 3, \cdots, k}
    \end{equation*}

    Prove $P(k + 1)$ is true. That is 
    \begin{align*}
        x_{n + 1} = 2\paren{-3}^{k + 1} + 10 \cdot 5^k. 
    \end{align*}

    By define 
    \begin{align*}
        x_{n + 1} &= 2x_k + 15x_{k - 1} \\
        &= 2 \bracks{2\paren{-3}^k + 10 \cdot 5^{k - 1}} + 15 \cdot 2 \bracks{2\paren{-3}^k + 10 \cdot 5^{k - 1}}\\
        &= \text{Chatgpt, help me expand it. }\\
        &= 2\paren{-3}^{k + 1} + 10 \cdot 5^{k - 1}
    \end{align*}

    for $k \geq 2$. 

    $P(k + 1)$ is true. So this $P(n)$ is true for all $n$ in $\N$. 
\end{proof}

\begin{example}
    Suppose $x_1 = 3, x_2 = 5$ and $x_n = 3x_{n - 1 } + 2x_{n - 2}, \forall n \geq 3$. 

    Prove $x_n < 4^n$ for all $n$ in $\N$. 
\end{example}
\begin{example}
    Suppose you have a $m \times n$ rectangularchocolate bar with $mn$ unit square. 

    Let $N = mn$. Prove $N - 1$ breaks requred to break the choc bar into unit square. 
\end{example}